\subsection{The extension of absolute values}

PROPOSITION 9. Let K be a field, L an algebraic extension of K and f an absolute value on K. Then f can be extended to an absolute value on L.

Suppose first that there exists a valuation \(v\) on \(\mathbf{K}\) with real values such that \(f(x) = e^{-v(x)}\) . There exists a valuation \(v'\) on \(\mathbf{L}\) whose restriction to \(\mathbf{K}\) is equivalent to \(v\) (§ 3, no. 3, Proposition 5). Then \(v'\) is of height 0 or 1 (no. 1, Corollary 2 to Proposition 1) and therefore may be assumed to have real values. The restriction of the mapping \(x \mapsto e^{-v'(x)}\) to \(\mathbf{K}\) is an absolute value equivalent to \(f\) and hence of the form \(f^s\) with \(s > 0\) (General Topology, Chapter IX, § 3, no. 2, Proposition 5). We conclude that

\[
x \mapsto e ^ {- v ^ {\prime} (x) / s}
\]

is an absolute value on \(\mathbf{L}\) extending \(f\) .

Suppose now that \(f\) is not ultrametric. Then \(K\) is identified with a subfield of \(\mathbf{C}\) such that \(f(x) = |x|^s\) where \(0 \leqslant s \leqslant 1\) (§ 6, no. 4, Theorem 2). As \(\mathbf{C}\) is algebraically closed, \(L\) is identified with a subfield of \(\mathbf{C}\) and the absolute value \(x \mapsto |x|^s\) extends \(f\) .

PROPOSITION 10. Let K be a field, f an absolute value on K such that K is complete and not discrete with respect to f and L an algebraic extension of K. Then f can be extended uniquely to an absolute value \(f'\) on L and, if L is of finite degree n, then

\[
f ^ {\prime} (x) = \left(f (\mathrm{N} _ {\mathrm{L/K}} (x))\right) ^ {1 / n}
\]

for all \(x \in \mathbf{L}\) .

The existence of \(f'\) follows from Proposition 9 and its uniqueness (over every finite sub-extension of \(L\) and therefore over the whole of \(L\) ) from Lemma 2 of §6, no. 4. Let \(f'\) be the unique extension of \(f\) to the algebraic closure of \(K\) and

suppose L is of finite degree n.~We know that \(\mathrm{N}_{\mathrm{L}/\mathrm{K}}(x)=\prod_{i=1}^{n} x_i\) , where each \(x_{i}\) is a conjugate of x over K (Algebra, Chapter VIII, § 12, no. 2, Proposition 4). In view of the uniqueness of \(f', f'(x_{i})=f'(x)\) for all i, whence the stated formula.

PROPOSITION 11. Let K be a field, f a non-ultrametric absolute value on K, \(\hat{K}\) the completion of K with respect to f, \(\hat{f}\) the continuous extension of f to \(\hat{K}\) and L a finite extension of K of degree n.

\begin{enumerate}
\def\labelenumi{(\alph{enumi})}
\item
  Let \(f'\) be an absolute value on \(\mathbf{L}\) extending \(f\) ; let \(\hat{\mathbf{L}}_{f'}\) denote the completion of \(\mathbf{L}\) with respect to \(f'\) and let \(\hat{\mathbf{K}}\) be identified with the closure of \(\mathbf{K}\) in \(\hat{\mathbf{L}}_{f'}\) ; then \([\hat{\mathbf{L}}_{f'}: \hat{\mathbf{K}}] \leqslant n\) .
\item
  The absolute values on L extending f are finite in number. If they are denoted by \(f_{1}^{\prime}, \ldots, f_{s}^{\prime}\) and the completion of L with respect to \(f_{i}^{\prime}\) by \(\hat{L}_{i}\) , the canonical mapping
\end{enumerate}

\(\hat{\mathbf{K}}\otimes_{\mathbf{K}}\mathbf{L}\rightarrow \prod_{i = 1}^{n}\hat{\mathbf{L}}_i\) is an isomorphism and

\[
\sum_ {i = 1} ^ {s} \left[ \hat {\mathrm{L}} _ {i}: \hat {\mathrm{K}} \right] = n.\tag{10}
\]

The proof is the same as that for the analogous assertions in Proposition 2 (no. 2). The references

§ 7, no. 2, Theorem 1; § 5, no. 2, Corollary to Proposition 4

should be replaced by the following

§ 7, no. 3, Theorem 2; Topological Vector Spaces, Chapter I,

§ 2, no. 3, Corollary 1 to Theorem 2.

Observe that two extensions of \(f\) to \(\mathbf{L}\) which define the same topology are equal (General Topology, Chapter IX, § 3, no. 2, Proposition 5). Finally, as \(f\) is not ultrametric, \(\mathbf{K}\) is of characteristic 0 and hence the Jacobson radical of \(\hat{\mathbf{K}} \otimes_{\mathbf{K}} \mathbf{L}\) is zero.

\textbf{Remarks}\par

\begin{enumerate}
\def\labelenumi{(\arabic{enumi})}
\item
  Proposition 11 (b) shows that every composite extension of \(\hat{\mathbf{K}}\) and \(\mathbf{L}\) over \(\mathbf{K}\) is isomorphic to one of the completions \(\hat{\mathbf{L}}_i\) and that these are composite extensions no two of which are isomorphic.
\item
  We know that the completions \(\hat{\mathbf{K}}\) and \(\hat{\mathbf{L}}_i\) are isomorphic to \(\mathbf{R}\) or \(\mathbf{C}\) (§ 6, no. 4, Theorem 2). If \(\hat{\mathbf{K}}\) is isomorphic to \(\mathbf{C}\) , so is \(\hat{\mathbf{L}}_i\) for all \(i\) and (10) shows that the number of extensions \(f_i'\) is equal to \(n\) . If \(\hat{\mathbf{K}}\) is isomorphic to \(\mathbf{R}\) (for example if \(\mathbf{K} = \mathbf{Q}\) ), let \(r_1\) (resp. \(r_2\) ) denote the number of indices \(i\) such that \(\hat{\mathbf{L}}_i\) is isomorphic to \(\mathbf{R}\) (resp. \(\mathbf{C}\) ); then (10) may be written:
\end{enumerate}

\[
r _ {1} + 2 r _ {2} = n.\tag{11}
\]

PROPOSITION 12. Let K be a field, f an absolute value on K, L a quasi-Galois extension of K and \(f'\) and \(f''\) two extensions of f to L. Then there exists a K-automorphism s of L such that \(f'' = f' \circ s\) .

If \(f\) is ultrametric, Corollary 1 to Proposition 7 (no. 6) shows that there exists a K-automorphism \(s\) of \(L\) such that \(f''\) and \(f' \circ s\) are equivalent absolute values; then there exists a real number \(a > 0\) such that \(f''(x) = (f'(s(x)))^a\) for all \(x \in L\) . If \(f\) is not improper, take \(x \in K^*\) such that \(f(x) \neq 1\) , which shows that \(a = 1\) . If \(f\) is improper, so are \(f'\) and \(f''\) (Corollary 2 to Proposition 1, no. 1) and \(s\) may be taken to be the identity automorphism.

If \(f\) is not ultrametric, there exist \(\mathbf{Q}\) -isomorphisms \(u', u''\) of \(\mathbf{L}\) onto subfields of \(\mathbf{C}\) and real exponents \(a' > 0\) , \(a'' > 0\) such that \(f'(x) = |u'(x)|^{a'}\) and

\[
f ^ {\prime \prime} (x) = | u ^ {\prime \prime} (x) | ^ {a ^ {\prime \prime}}
\]

for all \(x \in \mathbf{L}\) (§ 6, no. 4, Theorem 2). Taking \(x = 2\) , it is seen that \(a' = a''\) . The restrictions of \(u'\) and \(u''\) to \(\mathbf{K}\) extend by continuity to isomorphisms \(u_1\) and \(u_2\) of \(\hat{\mathbf{K}}\) onto \(\mathbf{R}\) (resp. \(\mathbf{C}\) ). Then \(u_2 \circ \bar{u}_1^{-1}\) is an automorphism of the valued field \(\mathbf{R}\) (resp. \(\mathbf{C}\) ) and is therefore the identity (resp. the identity or the automorphism \(c: \zeta \to \bar{\zeta}\) ). Replacing if need be \(u'\) by \(c \circ u'\) , it is seen that the restrictions of \(u'\) and \(u''\) to \(\mathbf{K}\) may be assumed to coincide. Identifying \(\mathbf{K}\) with a subfield of \(\mathbf{C}\) by means of this common restriction, \(u'\) and \(u''\) are K-isomorphisms of \(\mathbf{L}\) onto subfields of \(\mathbf{C}\) . As \(\mathbf{L}\) is a quasi-Galois extension of \(\mathbf{K}\) , there exists a K-automorphism \(s\) of \(\mathbf{L}\) such that \(u'' = u' \circ s\) ; since \(a' = a''\) , we deduce immediately that \(f'' = f' \circ s\) .

Remark (3). If \(\hat{\mathbf{K}}\) is isomorphic to \(\mathbf{R}\) , Proposition 12 shows that all the completions \(\hat{\mathbf{L}}_t\) of \(\mathbf{L}\) (in the notation of Proposition 11) are isomorphic to one another. Thus, with the notation of Remark 2 above, either \(r_1 = n\) and \(r_2 = 0\) , or \(r_1 = 0\) and \(2r_2 = n\) .

